\documentclass[11pt]{article}
\newcommand{\AssignmentNumber}{1}
\newcommand{\AssignmentTitle}{Build a Game with raylib}
\newcommand{\PostedDate}{Thursday, 17 September 2026}
\newcommand{\TeamSize}{exactly 3 students per group}
\newcommand{\DueDate}{Code + live presentation, on a day \textbf{TBA} during the week following the Final Examination period}
% Shared preamble for Computer Programming assignment sheets.
% Include with % Shared preamble for Computer Programming assignment sheets.
% Include with % Shared preamble for Computer Programming assignment sheets.
% Include with \input{common.tex} after \documentclass{article}.

\usepackage[a4paper,margin=2.2cm]{geometry}
\usepackage[T1]{fontenc}
\usepackage{lmodern}
\usepackage{microtype}
\usepackage{enumitem}
\usepackage{xcolor}
\usepackage{tcolorbox}
\tcbuselibrary{breakable,skins}
\usepackage{listings}
\usepackage{fancyhdr}
\usepackage{titlesec}
\usepackage{amsmath}
\usepackage{array}
\usepackage{booktabs}
\usepackage{hyperref}

\hypersetup{
  colorlinks=true,
  linkcolor=black,
  urlcolor=blue!60!black,
  pdfborder={0 0 0}
}

% ---------- Course identity (override per file if needed) ----------
\providecommand{\CourseName}{Computer Programming}
\providecommand{\CourseCode}{010153002}
\providecommand{\Semester}{Semester 1/2026}

% ---------- Colors ----------
\definecolor{accent}{HTML}{1F4E79}
\definecolor{accentlight}{HTML}{D9E2EC}
\definecolor{warn}{HTML}{B45309}
\definecolor{ok}{HTML}{166534}
\definecolor{codebg}{HTML}{F6F8FA}
\definecolor{codeframe}{HTML}{D0D7DE}
\definecolor{codekw}{HTML}{0550AE}
\definecolor{codestr}{HTML}{0A7B2C}
\definecolor{codecom}{HTML}{6E7781}

% ---------- Code listings ----------
\lstdefinestyle{c}{
  language=C,
  basicstyle=\ttfamily\small,
  keywordstyle=\color{codekw}\bfseries,
  stringstyle=\color{codestr},
  commentstyle=\color{codecom}\itshape,
  backgroundcolor=\color{codebg},
  frame=single,
  rulecolor=\color{codeframe},
  framerule=0.4pt,
  showstringspaces=false,
  breaklines=true,
  columns=fullflexible,
  keepspaces=true,
  tabsize=4,
  xleftmargin=8pt,
  xrightmargin=4pt,
  aboveskip=6pt,
  belowskip=6pt,
}
\lstset{style=c}

% ---------- Section styling ----------
\titleformat{\section}{\Large\bfseries\color{accent}}{\thesection.}{0.5em}{}
\titleformat{\subsection}{\large\bfseries\color{accent!85!black}}{\thesubsection}{0.5em}{}
\titlespacing*{\section}{0pt}{12pt}{6pt}
\titlespacing*{\subsection}{0pt}{8pt}{4pt}

% ---------- Page headers/footers ----------
\pagestyle{fancy}
\fancyhf{}
\fancyhead[L]{\small\CourseName\ifx\CourseCode\empty\else\ (\CourseCode)\fi}
\fancyhead[C]{\small\Semester}
\fancyhead[R]{\small\textbf{Assignment \AssignmentNumber:} \AssignmentTitle}
\fancyfoot[C]{\small\thepage}
\renewcommand{\headrulewidth}{0.4pt}

% ---------- Title block ----------
\newcommand{\assignmentheader}{%
  \begin{tcolorbox}[
    enhanced,
    colback=accentlight,
    colframe=accent,
    boxrule=0.5pt,
    arc=2pt,
    left=10pt, right=10pt, top=8pt, bottom=8pt
  ]
    {\Large\bfseries\color{accent} Programming Assignment \AssignmentNumber: \AssignmentTitle}\\[2pt]
    {\small \CourseName\ifx\CourseCode\empty\else\ (\CourseCode)\fi\quad\textbar\quad \Semester}\\[1pt]
    {\small \textbf{Posted:} \PostedDate\quad\textbar\quad \textbf{Team size:} \TeamSize}\\[1pt]
    {\small \textbf{Submission:} \DueDate}
  \end{tcolorbox}
  \vspace{2pt}
}

% ---------- Custom environments ----------
\newtcolorbox{summary}[1][Quick Reference]{
  enhanced,breakable,
  colback=accentlight!40,colframe=accent!70,
  title={\bfseries #1},
  fonttitle=\color{white},
  attach boxed title to top left={xshift=6pt,yshift=-3pt},
  boxed title style={colback=accent,boxrule=0pt,arc=2pt},
  arc=2pt,boxrule=0.4pt,
  top=12pt,
}

\newtcolorbox{rulebox}[1][The One Rule That Matters Most]{
  enhanced,breakable,
  colback=warn!8,colframe=warn,
  title={\bfseries #1},
  fonttitle=\color{white},
  attach boxed title to top left={xshift=6pt,yshift=-3pt},
  boxed title style={colback=warn,boxrule=0pt,arc=2pt},
  arc=2pt,boxrule=0.4pt,
  top=12pt,
}

\newtcolorbox{deliverables}[1][Deliverables]{
  enhanced,breakable,
  colback=ok!6,colframe=ok!70!black,
  title={\bfseries #1},
  fonttitle=\color{white},
  attach boxed title to top left={xshift=6pt,yshift=-3pt},
  boxed title style={colback=ok!70!black,boxrule=0pt,arc=2pt},
  arc=2pt,boxrule=0.4pt,
  top=12pt,
}
 after \documentclass{article}.

\usepackage[a4paper,margin=2.2cm]{geometry}
\usepackage[T1]{fontenc}
\usepackage{lmodern}
\usepackage{microtype}
\usepackage{enumitem}
\usepackage{xcolor}
\usepackage{tcolorbox}
\tcbuselibrary{breakable,skins}
\usepackage{listings}
\usepackage{fancyhdr}
\usepackage{titlesec}
\usepackage{amsmath}
\usepackage{array}
\usepackage{booktabs}
\usepackage{hyperref}

\hypersetup{
  colorlinks=true,
  linkcolor=black,
  urlcolor=blue!60!black,
  pdfborder={0 0 0}
}

% ---------- Course identity (override per file if needed) ----------
\providecommand{\CourseName}{Computer Programming}
\providecommand{\CourseCode}{010153002}
\providecommand{\Semester}{Semester 1/2026}

% ---------- Colors ----------
\definecolor{accent}{HTML}{1F4E79}
\definecolor{accentlight}{HTML}{D9E2EC}
\definecolor{warn}{HTML}{B45309}
\definecolor{ok}{HTML}{166534}
\definecolor{codebg}{HTML}{F6F8FA}
\definecolor{codeframe}{HTML}{D0D7DE}
\definecolor{codekw}{HTML}{0550AE}
\definecolor{codestr}{HTML}{0A7B2C}
\definecolor{codecom}{HTML}{6E7781}

% ---------- Code listings ----------
\lstdefinestyle{c}{
  language=C,
  basicstyle=\ttfamily\small,
  keywordstyle=\color{codekw}\bfseries,
  stringstyle=\color{codestr},
  commentstyle=\color{codecom}\itshape,
  backgroundcolor=\color{codebg},
  frame=single,
  rulecolor=\color{codeframe},
  framerule=0.4pt,
  showstringspaces=false,
  breaklines=true,
  columns=fullflexible,
  keepspaces=true,
  tabsize=4,
  xleftmargin=8pt,
  xrightmargin=4pt,
  aboveskip=6pt,
  belowskip=6pt,
}
\lstset{style=c}

% ---------- Section styling ----------
\titleformat{\section}{\Large\bfseries\color{accent}}{\thesection.}{0.5em}{}
\titleformat{\subsection}{\large\bfseries\color{accent!85!black}}{\thesubsection}{0.5em}{}
\titlespacing*{\section}{0pt}{12pt}{6pt}
\titlespacing*{\subsection}{0pt}{8pt}{4pt}

% ---------- Page headers/footers ----------
\pagestyle{fancy}
\fancyhf{}
\fancyhead[L]{\small\CourseName\ifx\CourseCode\empty\else\ (\CourseCode)\fi}
\fancyhead[C]{\small\Semester}
\fancyhead[R]{\small\textbf{Assignment \AssignmentNumber:} \AssignmentTitle}
\fancyfoot[C]{\small\thepage}
\renewcommand{\headrulewidth}{0.4pt}

% ---------- Title block ----------
\newcommand{\assignmentheader}{%
  \begin{tcolorbox}[
    enhanced,
    colback=accentlight,
    colframe=accent,
    boxrule=0.5pt,
    arc=2pt,
    left=10pt, right=10pt, top=8pt, bottom=8pt
  ]
    {\Large\bfseries\color{accent} Programming Assignment \AssignmentNumber: \AssignmentTitle}\\[2pt]
    {\small \CourseName\ifx\CourseCode\empty\else\ (\CourseCode)\fi\quad\textbar\quad \Semester}\\[1pt]
    {\small \textbf{Posted:} \PostedDate\quad\textbar\quad \textbf{Team size:} \TeamSize}\\[1pt]
    {\small \textbf{Submission:} \DueDate}
  \end{tcolorbox}
  \vspace{2pt}
}

% ---------- Custom environments ----------
\newtcolorbox{summary}[1][Quick Reference]{
  enhanced,breakable,
  colback=accentlight!40,colframe=accent!70,
  title={\bfseries #1},
  fonttitle=\color{white},
  attach boxed title to top left={xshift=6pt,yshift=-3pt},
  boxed title style={colback=accent,boxrule=0pt,arc=2pt},
  arc=2pt,boxrule=0.4pt,
  top=12pt,
}

\newtcolorbox{rulebox}[1][The One Rule That Matters Most]{
  enhanced,breakable,
  colback=warn!8,colframe=warn,
  title={\bfseries #1},
  fonttitle=\color{white},
  attach boxed title to top left={xshift=6pt,yshift=-3pt},
  boxed title style={colback=warn,boxrule=0pt,arc=2pt},
  arc=2pt,boxrule=0.4pt,
  top=12pt,
}

\newtcolorbox{deliverables}[1][Deliverables]{
  enhanced,breakable,
  colback=ok!6,colframe=ok!70!black,
  title={\bfseries #1},
  fonttitle=\color{white},
  attach boxed title to top left={xshift=6pt,yshift=-3pt},
  boxed title style={colback=ok!70!black,boxrule=0pt,arc=2pt},
  arc=2pt,boxrule=0.4pt,
  top=12pt,
}
 after \documentclass{article}.

\usepackage[a4paper,margin=2.2cm]{geometry}
\usepackage[T1]{fontenc}
\usepackage{lmodern}
\usepackage{microtype}
\usepackage{enumitem}
\usepackage{xcolor}
\usepackage{tcolorbox}
\tcbuselibrary{breakable,skins}
\usepackage{listings}
\usepackage{fancyhdr}
\usepackage{titlesec}
\usepackage{amsmath}
\usepackage{array}
\usepackage{booktabs}
\usepackage{hyperref}

\hypersetup{
  colorlinks=true,
  linkcolor=black,
  urlcolor=blue!60!black,
  pdfborder={0 0 0}
}

% ---------- Course identity (override per file if needed) ----------
\providecommand{\CourseName}{Computer Programming}
\providecommand{\CourseCode}{010153002}
\providecommand{\Semester}{Semester 1/2026}

% ---------- Colors ----------
\definecolor{accent}{HTML}{1F4E79}
\definecolor{accentlight}{HTML}{D9E2EC}
\definecolor{warn}{HTML}{B45309}
\definecolor{ok}{HTML}{166534}
\definecolor{codebg}{HTML}{F6F8FA}
\definecolor{codeframe}{HTML}{D0D7DE}
\definecolor{codekw}{HTML}{0550AE}
\definecolor{codestr}{HTML}{0A7B2C}
\definecolor{codecom}{HTML}{6E7781}

% ---------- Code listings ----------
\lstdefinestyle{c}{
  language=C,
  basicstyle=\ttfamily\small,
  keywordstyle=\color{codekw}\bfseries,
  stringstyle=\color{codestr},
  commentstyle=\color{codecom}\itshape,
  backgroundcolor=\color{codebg},
  frame=single,
  rulecolor=\color{codeframe},
  framerule=0.4pt,
  showstringspaces=false,
  breaklines=true,
  columns=fullflexible,
  keepspaces=true,
  tabsize=4,
  xleftmargin=8pt,
  xrightmargin=4pt,
  aboveskip=6pt,
  belowskip=6pt,
}
\lstset{style=c}

% ---------- Section styling ----------
\titleformat{\section}{\Large\bfseries\color{accent}}{\thesection.}{0.5em}{}
\titleformat{\subsection}{\large\bfseries\color{accent!85!black}}{\thesubsection}{0.5em}{}
\titlespacing*{\section}{0pt}{12pt}{6pt}
\titlespacing*{\subsection}{0pt}{8pt}{4pt}

% ---------- Page headers/footers ----------
\pagestyle{fancy}
\fancyhf{}
\fancyhead[L]{\small\CourseName\ifx\CourseCode\empty\else\ (\CourseCode)\fi}
\fancyhead[C]{\small\Semester}
\fancyhead[R]{\small\textbf{Assignment \AssignmentNumber:} \AssignmentTitle}
\fancyfoot[C]{\small\thepage}
\renewcommand{\headrulewidth}{0.4pt}

% ---------- Title block ----------
\newcommand{\assignmentheader}{%
  \begin{tcolorbox}[
    enhanced,
    colback=accentlight,
    colframe=accent,
    boxrule=0.5pt,
    arc=2pt,
    left=10pt, right=10pt, top=8pt, bottom=8pt
  ]
    {\Large\bfseries\color{accent} Programming Assignment \AssignmentNumber: \AssignmentTitle}\\[2pt]
    {\small \CourseName\ifx\CourseCode\empty\else\ (\CourseCode)\fi\quad\textbar\quad \Semester}\\[1pt]
    {\small \textbf{Posted:} \PostedDate\quad\textbar\quad \textbf{Team size:} \TeamSize}\\[1pt]
    {\small \textbf{Submission:} \DueDate}
  \end{tcolorbox}
  \vspace{2pt}
}

% ---------- Custom environments ----------
\newtcolorbox{summary}[1][Quick Reference]{
  enhanced,breakable,
  colback=accentlight!40,colframe=accent!70,
  title={\bfseries #1},
  fonttitle=\color{white},
  attach boxed title to top left={xshift=6pt,yshift=-3pt},
  boxed title style={colback=accent,boxrule=0pt,arc=2pt},
  arc=2pt,boxrule=0.4pt,
  top=12pt,
}

\newtcolorbox{rulebox}[1][The One Rule That Matters Most]{
  enhanced,breakable,
  colback=warn!8,colframe=warn,
  title={\bfseries #1},
  fonttitle=\color{white},
  attach boxed title to top left={xshift=6pt,yshift=-3pt},
  boxed title style={colback=warn,boxrule=0pt,arc=2pt},
  arc=2pt,boxrule=0.4pt,
  top=12pt,
}

\newtcolorbox{deliverables}[1][Deliverables]{
  enhanced,breakable,
  colback=ok!6,colframe=ok!70!black,
  title={\bfseries #1},
  fonttitle=\color{white},
  attach boxed title to top left={xshift=6pt,yshift=-3pt},
  boxed title style={colback=ok!70!black,boxrule=0pt,arc=2pt},
  arc=2pt,boxrule=0.4pt,
  top=12pt,
}


\begin{document}
\assignmentheader

\section*{The Mission}

Your team's mission is to design and build a \textbf{playable game} using
\href{https://www.raylib.com/}{\textbf{raylib}}, a simple C library for graphics, input, and audio.

\textbf{Your goal:} a complete, working game --- written in \textbf{plain C}, using only the
programming concepts covered in this course --- that your team can confidently explain, line by
line, in front of the class.

This is not just a coding exercise. It is a test of whether your team truly \textbf{understands}
every piece of code you submit.

\begin{rulebox}
\textbf{If you (or your teammates) cannot explain a piece of code, or it does not relate to what
we've studied in class, it earns you no score for that part --- no matter how impressive it looks.}

\medskip
Copy-pasted tutorial code, AI-generated boilerplate no one understands, and unrelated
libraries/features are all treated the same way: \textbf{unexplainable code = zero credit for
that section}, and it may reduce your whole-project score.
\end{rulebox}

\section*{Scope: What You May Use}

\begin{minipage}[t]{0.48\textwidth}
\textbf{Allowed:}
\begin{itemize}[leftmargin=*,nosep]
  \item \textbf{Plain C} (the course's C17/K\&R style) --- no C++ classes, STL, or \texttt{.cpp} files
  \item \textbf{raylib's C API only} --- \texttt{InitWindow}, \texttt{BeginDrawing}/\texttt{EndDrawing},
        \texttt{DrawText}, \texttt{DrawRectangle}, \texttt{IsKeyDown}, \texttt{LoadSound}, etc.
        (raylib is your ``screen and controller,'' not a shortcut around writing your own game logic)
  \item Every concept from \textbf{Lectures 1--10} of this course (see the checklist below)
\end{itemize}
\end{minipage}
\hfill
\begin{minipage}[t]{0.48\textwidth}
\textbf{Not allowed / flagged as ``unrelated code'':}
\begin{itemize}[leftmargin=*,nosep]
  \item C++ features (classes, \texttt{std::vector}, templates, \texttt{new}/\texttt{delete}, \ldots)
  \item External game engines, physics/AI libraries, or frameworks beyond raylib
  \item Multithreading, networking, or other techniques never covered in class
  \item Large chunks of code lifted from tutorials that your team cannot explain
\end{itemize}
\end{minipage}

\section*{Required Topic Coverage}

Your game's \textbf{logic} (not just its graphics) must clearly use the concepts below. You will
present a \textbf{coverage matrix} showing where each one lives in your code. Your team should aim
to \textbf{clearly demonstrate at least 8 of these 9} in working, explainable code. Using a topic
just to ``check a box'' without a real purpose will not count.

\begin{center}
\begin{tabular}{@{}clp{7.7cm}@{}}
\toprule
\# & Topic (Lecture) & Typical use in a game \\
\midrule
1 & Data types, variables, constants & Player stats, screen size constants \\
2 & Expressions \& operators & Score math, collision checks \\
3 & Selection (\texttt{if}/\texttt{switch}) & Game states, win/lose conditions \\
4 & Repetition (loops) & The main game loop, animations \\
5 & Functions, macros, recursion & Modular update/draw functions \\
6 & Arrays (1D/2D) + sorting/searching & Entities, grids, leaderboard \\
7 & Pointers \& dynamic memory (\texttt{malloc}/\texttt{free}) & Growing lists of bullets/enemies \\
8 & Strings & Player name, HUD text, messages \\
9 & \texttt{struct} / \texttt{union} / \texttt{enum} & Game objects, game-state enum \\
\bottomrule
\end{tabular}
\end{center}

\section*{Game Idea Starters}

Pick one, combine ideas, or bring your own (check with the instructor first):

\begin{itemize}[leftmargin=*,nosep]
  \item \textbf{Snake} --- array/grid of segments, struct for each segment, growing the snake with dynamic memory
  \item \textbf{Breakout / Brick Breaker} --- 2D array of bricks (struct per brick), collision via expressions
  \item \textbf{Space Shooter} --- dynamic array (pointers) of bullets/enemies, struct per entity, sorted high-score list
  \item \textbf{Memory / Card-Matching Game} --- 2D array shuffled with \texttt{rand()}, struct per card, string comparison
  \item \textbf{Tic-Tac-Toe / Connect Four} (vs.\ simple AI) --- 2D array board, recursion for the AI, enum for game state
  \item \textbf{Maze Game} --- 2D array maze, recursion or loops for movement/pathfinding
\end{itemize}

\section*{Getting Started: Minimal raylib Skeleton}

\begin{lstlisting}
#include "raylib.h"

int main(void) {
    const int screenWidth = 800;
    const int screenHeight = 450;

    InitWindow(screenWidth, screenHeight, "My Game");
    SetTargetFPS(60);

    while (!WindowShouldClose()) {
        // 1. Update game state (your functions, structs, arrays...)

        // 2. Draw
        BeginDrawing();
            ClearBackground(RAYWHITE);
            DrawText("Hello, raylib!", 190, 200, 20, LIGHTGRAY);
        EndDrawing();
    }

    CloseWindow();
    return 0;
}
\end{lstlisting}

Install raylib and see build instructions for your OS at \textbf{raylib.com} (Quickstart).

\section*{Weekly Progress Checkpoints}

Before the final code submission and presentation, each team checks in \textbf{three times},
about once a week, so problems get caught early instead of the night before the deadline.
Post each checkpoint to the course portal: a short status note (a few sentences) plus your
current source code (or a link/clip showing it running).

\begin{center}
\begin{tabular}{@{}lp{2.6cm}p{8.3cm}@{}}
\toprule
Checkpoint & Target date & What to show \\
\midrule
1 --- Concept \& Setup & Thu, 24 Sep 2026 & Team roster confirmed; game concept chosen; raylib
toolchain working (a window opens and draws something, even just a colored rectangle) \\
2 --- Core Loop & Thu, 1 Oct 2026 & The main game loop updates and draws every frame; at least
one \texttt{struct} and one array already in use for real game state \\
3 --- Feature-Complete Draft & Thu, 8 Oct 2026 & All core mechanics implemented and playable;
coverage matrix draft showing at least 5 of the 9 required topics in use \\
\bottomrule
\end{tabular}
\end{center}

These checkpoints are low-stakes and not separately scored, but a team that skips two or more of
them without a good reason should expect this to be reflected in the \textbf{Teamwork \&
Individual Contribution} portion of the final rubric below.

\begin{deliverables}
Submit a single ZIP file containing:
\begin{enumerate}[leftmargin=*,nosep]
  \item \textbf{Source Code} --- organized \texttt{.c}/\texttt{.h} files (no single 1000-line \texttt{main.c} --- use functions!)
  \item \textbf{Coverage Matrix} --- a short document (or slide) mapping each topic from Lectures 1--10 to the exact file/function where it's used
  \item \textbf{Design Report (PDF)} --- brief description of the game, a flowchart of the main game loop, and each team member's contribution
  \item \textbf{Compile Instructions} --- the exact command(s) to build your game (no Makefile knowledge required --- a plain \texttt{gcc ... -lraylib ...} command is fine)
\end{enumerate}
\end{deliverables}

\section*{The 15-Minute Presentation}

Each group presents \textbf{live, in front of the class}, for \textbf{15 minutes total}:

\begin{center}
\begin{tabular}{@{}lp{9cm}@{}}
\toprule
Time & Segment \\
\midrule
0--2 min & Game concept \& design overview (whole team) \\
2--8 min & Live gameplay demo \\
8--13 min & Code walkthrough --- each member explains \textbf{the part they personally wrote} \\
13--15 min & Instructor Q\&A --- any member may be asked about \textbf{any} part of the code \\
\bottomrule
\end{tabular}
\end{center}

\textbf{All 3 members must speak} and must each be able to answer questions about the whole
project, not only their own section.

\section*{Grading Rubric (100 points)}

\begin{center}
\begin{tabular}{@{}lcp{6.7cm}@{}}
\toprule
Category & Points & Focus \\
\midrule
Core Gameplay \& Functionality & 40 & Game runs, is playable, has clear win/lose logic \\
Concept Coverage \& Correct Use & 25 & Correct, purposeful use of $\geq$8 topics from the checklist \\
Code Quality \& Organization & 15 & Functions, meaningful names, no dead/unrelated code \\
Presentation \& Understanding & 10 & Every member explains their code clearly and fluently \\
Teamwork \& Individual Contribution & 10 & Balanced contribution across all 3 members \\
\bottomrule
\end{tabular}
\end{center}

\textbf{Reminder:} any code a team cannot explain, or that falls outside course topics, scores
\textbf{0} for that portion --- regardless of how the rest of the rubric is scored.

\section*{Submission}

\begin{itemize}[leftmargin=*,nosep]
  \item \textbf{What:} One ZIP per team (source code + coverage matrix + design report PDF)
  \item \textbf{Where:} Course portal
  \item \textbf{When:} Code submission and live presentation both happen on the same \textbf{TBA}
        day, sometime in the week following Final Examinations --- exact day and venue will be
        posted on the course portal once scheduled
\end{itemize}

Only one submission per team --- make sure all 3 names are listed in the design report.

\bigskip
\centerline{\textbf{Good luck! Build something you're proud of --- and something every teammate can explain with confidence.}}

\end{document}
